% Section 21 Header
\section{Section 21 — Miscellaneous Techniques}

\begin{frame}[c]{\empty}
    \begin{center}
        \Large\textbf{\secname}
    \end{center}
\end{frame}

% Slide 1: Concept
\begin{frame}{Miscellaneous Interview Techniques}
  Some interview questions don't fit into single patterns, but require a combination of general skills:
  \begin{itemize}
    \item \textbf{In-place Grid Simulation:} Modifying a matrix without using extra memory by using rows/columns headers as state markers (e.g. Set Matrix Zeroes).
    \item \textbf{Math \& Coordinate Transformations:} Mirroring, swapping, or transposing grids (e.g. Rotate Image).
    \item \textbf{Stack Evaluation:} Parsing arithmetic strings or expressions sequentially (e.g. Evaluate Reverse Polish Notation).
    \item \textbf{Bitmasking or Row/Column hashing:} Validating board boundaries quickly (e.g. Valid Sudoku).
  \end{itemize}
\end{frame}

% Problem 1: Set Matrix Zeroes
\begin{frame}[c]{\empty}
    \begin{center}
        \Large \textbf{Problem 1: \href{https://leetcode.com/problems/set-matrix-zeroes/}{Set Matrix Zeroes}}
    \end{center}
\end{frame}

% Problem 2: Valid Sudoku
\begin{frame}[c]{\empty}
    \begin{center}
        \Large \textbf{Problem 2: \href{https://leetcode.com/problems/valid-sudoku/}{Valid Sudoku}}
    \end{center}
\end{frame}

% Problem 3: Rotate Image
\begin{frame}[c]{\empty}
    \begin{center}
        \Large \textbf{Problem 3: \href{https://leetcode.com/problems/rotate-image/}{Rotate Image}}
    \end{center}
\end{frame}

% Problem 4: Evaluate Reverse Polish Notation
\begin{frame}[c]{\empty}
    \begin{center}
        \Large \textbf{Problem 4: \href{https://leetcode.com/problems/evaluate-reverse-polish-notation/}{Evaluate Reverse Polish Notation}}
    \end{center}
\end{frame}
